\chapter{Boolean Analysis — Invariance — Smoothing}
%%% generated by the blueprint pipeline from TCSlib.BooleanAnalysis.GaussianSpace.Invariance.Smoothing
\section{Overview}
Interval averaging gives quantitative smooth approximations of Lipschitz functions. Fourth-derivative expectation comparisons imply Lipschitz comparisons with an explicit constant.

\section{Declarations}

\begin{definition}[Translated interval average]
\lean{BooleanAnalysis.GaussianSpace.Invariance.intervalSmooth}
\leanok
For real $\eta,x$ and any real function $\psi$, the interval smoothing is $S_\eta\psi(x)=\eta^{-1}\int_0^\eta\psi(x+t)\,dt$, with value zero for $\eta=0$. The integral is oriented when $\eta<0$ and is assigned zero when its integrand is not integrable.
\end{definition}

\begin{theorem}[Derivative of an interval average]
\lean{BooleanAnalysis.GaussianSpace.Invariance.intervalSmooth_hasDerivAt}
\leanok
\uses{BooleanAnalysis.GaussianSpace.Invariance.intervalSmooth}
\difficulty{4}
For $\eta>0$ and a real function $\psi$, write $S_\eta\psi(x)=\eta^{-1}\int_0^\eta\psi(x+t)\,dt$. If $\psi$ is continuous, then $S_\eta\psi$ is differentiable at every real $x$ with derivative $(\psi(x+\eta)-\psi(x))/\eta$.
\end{theorem}

\begin{theorem}[One additional derivative from averaging]
\lean{BooleanAnalysis.GaussianSpace.Invariance.contDiff_intervalSmooth}
\leanok
\uses{BooleanAnalysis.GaussianSpace.Invariance.intervalSmooth, BooleanAnalysis.GaussianSpace.Invariance.intervalSmooth_hasDerivAt}
\difficulty{4}
For $\eta>0$ and a real function $\psi$, write $S_\eta\psi(x)=\eta^{-1}\int_0^\eta\psi(x+t)\,dt$. For each integer $m\ge0$, if $\psi$ is $C^m$ on $\mathbb R$, then $S_\eta\psi$ is $C^{m+1}$ on $\mathbb R$.
\end{theorem}

\begin{theorem}[Higher derivatives of an interval average]
\lean{BooleanAnalysis.GaussianSpace.Invariance.iteratedDeriv_intervalSmooth}
\leanok
\uses{BooleanAnalysis.GaussianSpace.Invariance.intervalSmooth, BooleanAnalysis.GaussianSpace.Invariance.intervalSmooth_hasDerivAt}
\difficulty{4}
For $\eta>0$ and a real function $\psi$, write $S_\eta\psi(x)=\eta^{-1}\int_0^\eta\psi(x+t)\,dt$. If $m\ge0$ is an integer and $\psi\in C^m(\mathbb R)$, then $(S_\eta\psi)^{(m+1)}(x)=(\psi^{(m)}(x+\eta)-\psi^{(m)}(x))/\eta$ for every real $x$.
\end{theorem}

\begin{theorem}[Derivative bound inherited from Lipschitz continuity]
\lean{BooleanAnalysis.GaussianSpace.Invariance.abs_deriv_intervalSmooth_le}
\leanok
\uses{BooleanAnalysis.GaussianSpace.Invariance.intervalSmooth, BooleanAnalysis.GaussianSpace.Invariance.intervalSmooth_hasDerivAt}
\difficulty{3}
For $\eta>0$ and a real function $\psi$, write $S_\eta\psi(x)=\eta^{-1}\int_0^\eta\psi(x+t)\,dt$. If $\psi$ is $c$-Lipschitz for $c\ge0$, then $|(S_\eta\psi)\prime(x)|\le c$ for every real $x$.
\end{theorem}

\begin{theorem}[Interval averaging preserves the Lipschitz constant]
\lean{BooleanAnalysis.GaussianSpace.Invariance.lipschitz_intervalSmooth}
\leanok
\uses{BooleanAnalysis.GaussianSpace.Invariance.abs_deriv_intervalSmooth_le, BooleanAnalysis.GaussianSpace.Invariance.intervalSmooth, BooleanAnalysis.GaussianSpace.Invariance.intervalSmooth_hasDerivAt}
\difficulty{3}
For $\eta>0$ and a real function $\psi$, write $S_\eta\psi(x)=\eta^{-1}\int_0^\eta\psi(x+t)\,dt$. If $\psi\colon\mathbb R\to\mathbb R$ is $c$-Lipschitz with $c\ge0$, then $S_\eta\psi$ is also $c$-Lipschitz.
\end{theorem}

\begin{theorem}[Uniform approximation by interval averaging]
\lean{BooleanAnalysis.GaussianSpace.Invariance.abs_intervalSmooth_sub_le}
\leanok
\uses{BooleanAnalysis.GaussianSpace.Invariance.intervalSmooth}
\difficulty{4}
For $\eta>0$ and a real function $\psi$, write $S_\eta\psi(x)=\eta^{-1}\int_0^\eta\psi(x+t)\,dt$. If $\psi$ is $c$-Lipschitz for $c\ge0$, then $|S_\eta\psi(x)-\psi(x)|\le c\eta$ for every real $x$.
\end{theorem}

\begin{theorem}[Bounds for repeated interval averaging]
\lean{BooleanAnalysis.GaussianSpace.Invariance.intervalSmooth_iterate_bounds}
\leanok
\uses{BooleanAnalysis.GaussianSpace.Invariance.abs_deriv_intervalSmooth_le, BooleanAnalysis.GaussianSpace.Invariance.abs_intervalSmooth_sub_le, BooleanAnalysis.GaussianSpace.Invariance.contDiff_intervalSmooth, BooleanAnalysis.GaussianSpace.Invariance.intervalSmooth, BooleanAnalysis.GaussianSpace.Invariance.iteratedDeriv_intervalSmooth, BooleanAnalysis.GaussianSpace.Invariance.lipschitz_intervalSmooth}
\difficulty{5}
For $\eta>0$ and a real function $\psi$, write $S_\eta\psi(x)=\eta^{-1}\int_0^\eta\psi(x+t)\,dt$. If $\psi$ is $c$-Lipschitz with $c\ge0$ and $m\ge0$ is an integer, put $\phi=S_\eta^{m+1}\psi$, where the exponent denotes repeated application of $S_\eta$. Then $\phi$ is $C^{m+1}$ and $c$-Lipschitz, and for every real $x$, $|\phi(x)-\psi(x)|\le(m+1)c\eta$ and $|\phi^{(m+1)}(x)|\le c(2/\eta)^m$.
\end{theorem}

\begin{theorem}[Quantitative fourth-order smooth approximation]
\lean{BooleanAnalysis.GaussianSpace.Invariance.exists_contDiff_four_approx}
\leanok
\uses{BooleanAnalysis.GaussianSpace.Invariance.intervalSmooth, BooleanAnalysis.GaussianSpace.Invariance.intervalSmooth_iterate_bounds}
\difficulty{4}
For every $\eta>0$, $c\ge0$, and $c$-Lipschitz function $\psi\colon\mathbb R\to\mathbb R$, there exists $\phi\in C^4(\mathbb R)$ such that $|\phi(x)-\psi(x)|\le4c\eta$ and $|\phi^{(4)}(x)|\le8c/\eta^3$ for all real $x$.
\end{theorem}

\begin{theorem}[Expectation stability under uniform approximation]
\lean{BooleanAnalysis.GaussianSpace.Invariance.abs_integral_sub_le_uniform}
\leanok
\difficulty{4}
Let $f,g\colon\Omega\to\mathbb R$ be functions on a probability space $(\Omega,\mu)$, with $f$ integrable and $g$ measurable up to a null set. If a real constant $D$ satisfies $|f(\omega)-g(\omega)|\le D$ for every $\omega$, then $|\int f\,d\mu-\int g\,d\mu|\le D$.
\end{theorem}

\begin{theorem}[Lipschitz comparison at any smoothing scale]
\lean{BooleanAnalysis.GaussianSpace.Invariance.lipschitz_expectation_bound_at_scale}
\leanok
\uses{BooleanAnalysis.GaussianSpace.Invariance.abs_integral_sub_le_uniform, BooleanAnalysis.GaussianSpace.Invariance.exists_contDiff_four_approx, BooleanAnalysis.GaussianSpace.integrable_test_of_memLp_four}
\difficulty{4}
Let $U\in L^4(\mu)$ and $V\in L^4(\nu)$ be real inputs on probability spaces, and let $B\ge0$. Suppose that for every $\phi\in C^4(\mathbb R)$ and $C\ge0$ bounding $|\phi^{(4)}|$ everywhere, $|\mathbb E_\mu\phi(U)-\mathbb E_\nu\phi(V)|\le CB$. For every $c\ge0$, $c$-Lipschitz function $\psi\colon\mathbb R\to\mathbb R$, and $\eta>0$, $|\mathbb E_\mu\psi(U)-\mathbb E_\nu\psi(V)|\le8c\eta+(8c/\eta^3)B$.
\end{theorem}

\begin{theorem}[Lipschitz comparison from fourth-order tests]
\lean{BooleanAnalysis.GaussianSpace.Invariance.lipschitz_expectation_bound_of_fourth_derivative}
\leanok
\uses{BooleanAnalysis.GaussianSpace.Invariance.lipschitz_expectation_bound_at_scale}
\difficulty{5}
Let $U\in L^4(\mu)$ and $V\in L^4(\nu)$ be real inputs on probability spaces, and let $B\ge0$. Suppose that for every $\phi\in C^4(\mathbb R)$ and $C\ge0$ bounding $|\phi^{(4)}|$ everywhere, $|\mathbb E_\mu\phi(U)-\mathbb E_\nu\phi(V)|\le CB$. For every $c\ge0$ and $c$-Lipschitz function $\psi\colon\mathbb R\to\mathbb R$, $|\mathbb E_\mu\psi(U)-\mathbb E_\nu\psi(V)|\le16cB^{1/4}$, including $B=0$.
\end{theorem}
